% The print edition uses prerex to draw this as a page-sized dependency graph.
% LaTeXML has no prerex binding, so the web edition exposes the same guidance as
% a semantic list whose part references remain navigable.
\begin{quote}
PDF판의 선수관계 도표를 웹에서 읽기 쉬운 순서로 풀어 쓴 것입니다.
\begin{itemize}
  \item 출발점: 추상대수학 입문과 기본 위상수학.
  \item 선형대수학은 표현론, 양자 알고리즘, 미분기하학의 바탕입니다.
  \item 군론과 갈루아 이론은 대수적 정수론으로 이어집니다.
  \item 기본 위상수학과 군작용은 대수적 위상수학으로 이어집니다.
  \item 범주론은 표현론, 대수적 위상수학, 대수기하학을 잇는 공통 언어입니다.
  \item 미적분학과 위상수학은 복소해석학, 측도론, 미분기하학의 바탕입니다.
\end{itemize}
각 부와 장으로 가는 링크는 이 페이지의 목차에서 바로 열 수 있습니다.
\end{quote}
